% !TeX root = ../main.tex
\section{Mô Hình Toán Học Hình Thức (Chế Độ 2 Người)}
\label{app:formal}

Phụ lục này đặc tả Tactical Formation bằng toán học, \emph{độc lập với mọi
cách cài đặt phần mềm}, bổ sung cho phần trình bày chính bằng hình ảnh bàn
cờ thật ở Mục~\ref{sec:rules} (đủ để nắm luật chơi). Trọng tâm đặt
vào cốt lõi của cờ: quân đi theo \textbf{hướng} nào và \textbf{ăn quân}
ra sao. Mô hình dựng trên ba trụ cột---\textbf{trạng thái}, \textbf{hành
động}, \textbf{phần thưởng}---và một tiên đề xuyên suốt: \emph{bàn cờ
luôn có đúng 28 quân; ăn quân là đổi phe, không bao giờ loại bỏ}. Phần
mở rộng cho chế độ 3--4 người được nêu ở Mục~\ref{sec:variants}.

\subsection{Bàn cờ và 28 quân}
\label{ssec:board}

Bàn cờ chế độ 2 người gồm 41 ô trên hai lớp:
\begin{align}
\mathcal{B}_{\mathsf V}
  &= \{(i,j)\in\mathbb Z^2 : 0\le i,j\le 4\}
  && \text{(25 ô đỉnh)}, \label{eq:vert}\\
\mathcal{B}_{\mathsf C}
  &= \{(i{+}\tfrac12,\,j{+}\tfrac12) : 0\le i,j\le 3\}
  && \text{(16 ô tâm)}, \label{eq:cent}\\
\mathcal{B}
  &= \mathcal{B}_{\mathsf V}\sqcup\mathcal{B}_{\mathsf C},
  && |\mathcal{B}|=41. \label{eq:board}
\end{align}
Mỗi ô tâm nằm chính giữa bốn ô đỉnh bao quanh. Hàm loại
$\kappa(x)\in\{\mathsf V,\mathsf C\}$ cho biết ô $x$ là đỉnh hay tâm.
Hai phe là $\mathcal C=\{\mathrm Y,\mathrm G\}$ (Vàng, Xanh); đối thủ của
$c$ viết là $\bar c$; ô trống ký hiệu $\varnothing$. \textbf{Xanh đi
trước.}

Trò chơi có đúng \textbf{28 quân}; mỗi phe xuất phát với 14 quân
($=10$ quân đỉnh $+\,4$ quân tâm). Mỗi quân $p$ mang ba thuộc tính:
\begin{itemize}[noitemsep,topsep=2pt]
  \item \textbf{loại} $\kappa(p)\in\{\mathsf V,\mathsf C\}$---\emph{bất
        biến suốt ván}: quân đỉnh sống cả đời trên ô đỉnh, quân tâm trên
        ô tâm;
  \item \textbf{vị trí} $x_p$---đổi khi quân nhảy;
  \item \textbf{phe} $c_p\in\{\mathrm Y,\mathrm G\}$---đổi khi quân bị ăn.
\end{itemize}

\begin{quote}
\textbf{Định luật Bảo toàn 28 quân.} Không quân nào được thêm vào hay
bị loại khỏi bàn; ăn quân là \emph{đổi phe} một quân có sẵn. Do đó tại
mọi trạng thái
\begin{equation}
N_{\mathrm Y}(s)+N_{\mathrm G}(s)=28,
\label{eq:conservation}
\end{equation}
trong đó $N_c(s)$ là số quân phe $c$.
\end{quote}

Cần nói rõ địa vị của khẳng định này: đây là một \emph{bất biến hệ quả
trực tiếp của định nghĩa}---vì bắt quân được định nghĩa là \emph{lật màu}
một quân có sẵn (không thêm, không loại), tổng số quân hiển nhiên không
đổi---chứ không phải một định lý sâu cần chứng minh công phu. Giá trị của
nó nằm ở \emph{hệ quả đối với lối chơi}: ván cờ không bao giờ suy biến
thành tàn cuộc thưa quân, toàn bộ 28 quân luôn là tài nguyên chiến thuật
sống.

Bất biến~\eqref{eq:conservation} phân biệt Tactical Formation với cả Cờ
Đam (ăn $=$ loại) lẫn Reversi (số quân tăng theo lượt đặt), và xếp nó vào
cùng nhóm với Cờ gánh và Ming Mang, hai trò chơi cũng bảo toàn tổng số quân
(Mục~\ref{sec:related}). Vị trí xuất
phát (trạng thái $0$) gồm hai khối quân đối xứng:
\begin{align}
P^{\mathrm Y}_0 &= \{(i,0)\}_{i=0}^{4}\cup\{(i{+}\tfrac12,\tfrac12)\}_{i=0}^{3}\cup\{(i,1)\}_{i=0}^{4},
\label{eq:starty}\\
P^{\mathrm G}_0 &= \{(i,4)\}_{i=0}^{4}\cup\{(i{+}\tfrac12,\tfrac72)\}_{i=0}^{3}\cup\{(i,3)\}_{i=0}^{4}.
\label{eq:startg}
\end{align}

\subsection{Trạng thái}
\label{ssec:state}

Trạng thái trò chơi là cặp
\begin{equation}
\boxed{\;s=(b,\mathrm{turn}),\qquad
b:\mathcal B\to\{\mathrm Y,\mathrm G,\varnothing\},\qquad
\mathrm{turn}\in\{\mathrm Y,\mathrm G\}.\;}
\label{eq:state}
\end{equation}
$b$ là \emph{cấu hình bàn} (ô nào mang màu gì), $\mathrm{turn}$ là phe
đến lượt. Biết vị trí và phe của cả 28 quân cộng với ai đang đi là biết
toàn bộ trạng thái---trò chơi \emph{thông tin hoàn hảo}, không có thông
tin ẩn. Trạng thái xuất phát là $s_0=(b_0,\mathrm G)$ với $b_0$ đặt theo
\eqref{eq:starty}--\eqref{eq:startg}. Hàm đếm quân $N_c(s)$ thoả
$N_{\mathrm Y}(s_0)=N_{\mathrm G}(s_0)=14$ và bất biến~\eqref{eq:conservation}.

\subsection{Hướng đi và cú nhảy}
\label{ssec:action}

Tập hướng đi phụ thuộc loại quân:
\begin{equation}
D_{\mathsf V} = \underbrace{\{(\pm1,0),(0,\pm1)\}}_{\text{4 ngang--dọc, bước }1}
  \cup \underbrace{\{(\pm\tfrac12,\pm\tfrac12)\}}_{\text{4 chéo, bước }\frac12},
\qquad
D_{\mathsf C} = \{(\pm\tfrac12,\pm\tfrac12)\}.
\label{eq:dirs}
\end{equation}
Quân đỉnh có 8 hướng; quân tâm chỉ 4 hướng chéo (đi ngang--dọc một ô sẽ
rơi khỏi lớp tâm). Quân tại ô $x$ đi theo $d\in D_{\kappa(x)}$ bằng cách
nhảy \emph{qua} ô kề $x+d$ và \emph{đáp} xuống $x+2d$. Cú nhảy hợp lệ khi
và chỉ khi
\begin{equation}
\underbrace{x{+}d,\ x{+}2d\in\mathcal B}_{\text{không ra ngoài bàn}},
\qquad
\underbrace{b(x{+}d)\neq\varnothing}_{\text{bệ nhảy có quân (phe nào cũng được)}},
\qquad
\underbrace{b(x{+}2d)=\varnothing}_{\text{ô đáp trống}}.
\label{eq:jump}
\end{equation}
Vì $2d$ là vectơ nguyên nên ô đáp tự động cùng loại với $x$ (đỉnh về
đỉnh, tâm về tâm). \textbf{Quân làm bệ nhảy đứng nguyên}---không đổi chỗ,
không đổi phe. Tập ô đáp hợp lệ và tập nước di chuyển là
\begin{equation}
\begin{aligned}
\mathrm{Jumps}(x,b)&=\{x{+}2d : d\in D_{\kappa(x)},\ \eqref{eq:jump}\ \text{thoả}\},\\
\mathcal A_{\mathrm{move}}(s)&=\{(x,x') : b(x)=\mathrm{turn},\ x'\in\mathrm{Jumps}(x,b)\}.
\end{aligned}
\label{eq:moves}
\end{equation}
Bệ nhảy của quân đỉnh đi chéo là một ô \emph{tâm}; bệ của quân tâm là một
ô \emph{đỉnh}---hai lớp đan nhau qua đường chéo (xem minh hoạ ở
Hình~\ref{fig:jumps}).

\subsection{Cách ăn quân---cốt lõi của cờ}
\label{ssec:capture}

Ý tưởng trong một câu: \emph{đáp xuống đâu, nhìn quanh đó}---từ ô đáp,
theo từng hướng, nếu thấy một dãy quân địch liền nhau mà ngay sau dãy là
một quân mình, thì cả dãy bị \emph{kẹp} và \emph{đổi phe} theo mình.
Hướng ăn dùng cùng tập hướng với hướng đi: quân đỉnh xét
$\Delta_{\mathsf V}=D_{\mathsf V}$ (8 hướng), quân tâm xét
$\Delta_{\mathsf C}=D_{\mathsf C}$ (4 hướng chéo). \emph{Tia} theo hướng
$\delta$ từ ô đáp $x'$ là dãy ô $x'+\delta,\ x'+2\delta,\ x'+3\delta,\dots$;
tia ngang--dọc chỉ đi qua ô đỉnh, còn tia chéo đi xen kẽ tâm--đỉnh--tâm
nên một cú kẹp chéo có thể ăn lẫn cả quân đỉnh và quân tâm.

Gọi $b'$ là bàn ngay sau cú nhảy (đã nhảy, chưa lật) và $c=\mathrm{turn}$.
Hướng $\delta$ \textbf{ăn được} khi tồn tại số nguyên $m\ge 1$ sao cho
\begin{equation}
\underbrace{b'(x'+k\delta)=\bar c,\quad k=1,\dots,m}_{m\ \text{quân địch đứng liền nhau}},
\qquad
\underbrace{b'(x'+(m{+}1)\delta)=c}_{\text{ngay sau dãy là quân mình}}.
\label{eq:clamp}
\end{equation}
Khi đó tập lật của hướng $\delta$ là $F_\delta=\{x'+k\delta : 1\le k\le m\}$.
Hướng $\delta$ \emph{không} ăn được (quy ước $F_\delta=\varnothing$) nếu
tia---trước khi gặp quân mình---gặp ô trống, ra ngoài bàn, hoặc gặp ngay
một quân mình (chưa có quân địch nào ở giữa). Tập hướng ăn được của nước
$(x,x')$ là $\Delta(s,x,x')=\{\delta : F_\delta\neq\varnothing\}$.

\paragraph{Hành động đầy đủ.} Một hành động là bộ ba
\begin{equation}
\boxed{\;a=(x,\ x',\ \delta)\;}
\qquad
\delta\in
\begin{cases}
\Delta(s,x,x') & \text{nếu có hướng ăn được},\\
\{\varnothing\} & \text{nếu không (không lật gì)},
\end{cases}
\label{eq:action}
\end{equation}
với $(x,x')\in\mathcal A_{\mathrm{move}}(s)$. \textbf{Luật:} có kẹp thì
\emph{phải ăn}, và người đi chọn \emph{đúng một hướng}---các hướng ăn
được khác không bị lật. Mọi $F_\delta$ tính được ngay từ bàn-sau-cú-nhảy
nên người đi luôn thấy trước mỗi lựa chọn ăn được bao nhiêu quân.
\paragraph{Ví dụ.} \emph{Quy ước ký hiệu trong ví dụ:} chữ HOA
$\mathrm Y,\mathrm G$ chỉ quân \emph{đỉnh}, chữ thường $\mathrm y,\mathrm g$
chỉ quân \emph{tâm} (cùng màu Vàng/Xanh); đây chỉ là cách viết gọn để nhìn
thấy lớp của từng quân---hàm bàn $b$ ở~\eqref{eq:state} vẫn chỉ trả màu, còn
lớp là thuộc tính $\kappa(x)$ của ô. Quân Xanh ($c=\mathrm G$) vừa đáp xuống
$x'=(1,2)$ với cấu hình quanh đó: $b'(2,2)=\mathrm Y$, $b'(3,2)=\mathrm G$ trên tia
ngang; và $b'(1.5,2.5)=\mathrm y$, $b'(2,3)=\mathrm Y$,
$b'(2.5,3.5)=\mathrm y$, $b'(3,4)=\mathrm G$ trên tia chéo. Hai hướng ăn
được:
\begin{align}
\delta_1=(1,0):\quad & F_{\delta_1}=\{(2,2)\}, && |{\cdot}|=1, \label{eq:ex1}\\
\delta_2=(\tfrac12,\tfrac12):\quad & F_{\delta_2}=\{(1.5,2.5),(2,3),(2.5,3.5)\}, && |{\cdot}|=3. \label{eq:ex2}
\end{align}
Tia chéo~\eqref{eq:ex2} ăn lẫn cả quân tâm và quân đỉnh (2 tâm $+$ 1
đỉnh), minh hoạ việc tia chéo đi xen kẽ hai lớp. Người đi chọn đúng một
hướng; chọn $\delta_2$ thì ba quân Vàng trên tia chéo đổi sang Xanh, còn
$Y@(2,2)$ của tia ngang \emph{vẫn Vàng}. Tình huống này được minh hoạ
trực quan ở Hình~\ref{fig:capture} và Bảng~\ref{tab:capture_rays}.

\subsection{Phần thưởng}
\label{ssec:reward}

Cỡ của hành động là số quân nó ăn, $|a|:=|F_\delta|$ (bằng $0$ nếu
$\delta=\varnothing$). Phần thưởng từng bước là tổng-zero:
\begin{equation}
\boxed{\;r_c(s,a)=+|a|,\qquad r_{\bar c}(s,a)=-|a|\;}
\qquad (c=\mathrm{turn}).
\label{eq:reward}
\end{equation}
Bốn hệ quả ``kế toán''---tất cả đều là hệ quả của \emph{ăn $=$ đổi phe}:
\begin{enumerate}[label=(\arabic*),noitemsep,topsep=2pt]
  \item \textbf{Tổng-zero từng bước:} $r_{\mathrm Y}(s,a)+r_{\mathrm G}(s,a)=0$.
  \item \textbf{Ăn $k$ quân đổi chênh lệch quân số $2k$:} mình $+k$,
        địch $-k$ cùng lúc.
  \item \textbf{Tổng luôn là 28:} bất biến~\eqref{eq:conservation} ở mọi
        trạng thái.
  \item \textbf{Bảng điểm chính là bàn cờ:} dọc ván
        $s_0,a_0,s_1,\dots,s_t$,
        \begin{equation}
        N_c(s_t)=14+\sum_{k=0}^{t-1} r_c(s_k,a_k);
        \label{eq:scoreboard}
        \end{equation}
        đếm quân mỗi phe tại bất kỳ lúc nào là đọc được tổng phần thưởng
        tích luỹ của phe đó.
\end{enumerate}
Khi ván kết thúc, phe thắng nhận $+W$, phe thua $-W$ với $W>0$ là hằng số
trừu tượng (giá trị cụ thể do từng cài đặt chọn); tổng phần thưởng cả ván
của hai phe luôn bằng $0$.

\subsection{Hàm chuyển $T$}
\label{ssec:transition}

Cho $s=(b,\mathrm{turn})$, $a=(x,x',\delta)$, $c=\mathrm{turn}$. Một
nước đi diễn ra theo ba bước.

\noindent\textbf{Bước 1---nhảy.} Quân rời $x$, đáp $x'$:
\begin{equation}
b'(y)=\begin{cases}\varnothing & y=x,\\ c & y=x',\\ b(y) & \text{ô khác}.\end{cases}
\label{eq:t1}
\end{equation}
\textbf{Bước 2---lật theo hướng đã chọn.} Mọi quân trong $F_\delta$ đổi
sang màu $c$ (nếu $\delta=\varnothing$ thì không lật gì):
\begin{equation}
b''(y)=\begin{cases}c & y\in F_\delta,\\ b'(y) & \text{ô khác}.\end{cases}
\label{eq:t2}
\end{equation}
\textbf{Bước 3---đổi lượt.}
\begin{equation}
\boxed{\;T(s,a)=(b'',\,\bar c).\;}
\label{eq:t3}
\end{equation}
$T$ là hàm \emph{tất định}---không có yếu tố ngẫu nhiên nào trong cờ này.

\subsection{Điều kiện kết thúc}
\label{ssec:terminal}

Phe $c$ \emph{sống} ở trạng thái $s$ khi còn ít nhất một quân còn nước
nhảy:
\begin{equation}
\mathrm{alive}(c,s)\iff \exists\,x:\ b(x)=c\ \wedge\ \mathrm{Jumps}(x,b)\neq\varnothing.
\label{eq:alive}
\end{equation}
\begin{equation}
\boxed{\;\text{Đến lượt phe nào mà phe đó không còn sống}
\ \Rightarrow\ \text{phe đó thua ngay, phe kia thắng.}\;}
\label{eq:terminal}
\end{equation}
Hai cách thua tương ứng (Hình~\ref{fig:endgame}): \textbf{hết quân} (cả
14 quân gốc lẫn quân từng ăn được đều đã bị đổi phe) hoặc \textbf{bị phong
toả} (còn quân nhưng không quân nào nhảy được). Một phe tạm hết nước khi
\emph{chưa} tới lượt mình thì chưa thua---chỉ xét đúng lúc đến lượt. Chế
độ 2 người \emph{không có hoà}.

\subsection{Chặn trên độ phức tạp}
\label{ssec:complexity}

Định luật Bảo toàn~\eqref{eq:conservation} cho phép chặn kích thước không
gian trạng thái khá gọn. Loại quân bất biến nên luôn có đúng 20 quân đỉnh
(trên 25 ô đỉnh) và 8 quân tâm (trên 16 ô tâm); mỗi quân mang một trong
hai màu. Bỏ qua tính khả đạt, số cấu hình bàn bị chặn bởi
\begin{equation}
|\mathcal S| \;\le\; 2\cdot
\underbrace{\binom{25}{20}2^{20}}_{\text{lớp đỉnh}}\cdot
\underbrace{\binom{16}{8}\,2^{8}}_{\text{lớp tâm}}
\;\approx\; 3.7\times10^{17},
\label{eq:statespace}
\end{equation}
(thừa số $2$ cho lượt đi). Đây là chặn trên lỏng---nhỏ hơn nhiều bậc so
với Cờ Đam ($\sim5\times10^{20}$, đã giải~\citep{schaeffer2007checkers})
---cho thấy chế độ 2 người nằm trong tầm tìm kiếm heuristic mà không cần
học sâu.

Về \emph{branching factor}, một hành động là bộ ba $(x,x',\delta)$. Vì ăn
quân là đổi phe, một phe có thể nắm gần như toàn bộ 28 quân, nên không thể
chặn theo 14 quân ban đầu. Nhưng \emph{loại} quân thì bất biến: toàn bàn
luôn có đúng 20 quân đỉnh và 8 quân tâm. Quân đỉnh có $\le 8$ ô đáp, và từ
mỗi ô đáp có $\le 8$ tia kẹp để chọn một (nếu không có tia nào thì chỉ có
một hành động); quân tâm có $\le 4$ ô đáp và $\le 4$ tia. Do đó số hành
động ở một trạng thái bị chặn lỏng bởi
\begin{equation}
|\mathcal A(s)| \;\le\; 20\cdot 8\cdot 8 + 8\cdot 4\cdot 4 \;=\; 1408 .
\label{eq:branchbound}
\end{equation}
Ràng buộc cú nhảy~\eqref{eq:jump} (cần bệ \emph{và} ô đáp trống) khiến
branching factor thực tế nhỏ hơn nhiều; giá trị đo đạc được báo cáo ở
Mục~\ref{sec:eval}.

\subsection{Trò chơi trong một dòng}
\label{ssec:onelinerule}

Gộp lại, mô hình hoàn chỉnh là bộ
\begin{equation}
\boxed{\;\mathcal G=\big(\mathcal S,\ \mathcal A,\ T,\ r,\ R,\ s_0\big)\;}
\label{eq:game}
\end{equation}
---một trò chơi \emph{hai người, tổng-zero, thông tin hoàn hảo, lượt thay
phiên, tất định}: trạng thái~\eqref{eq:state}, hành động $=$ nhảy $+$
hướng ăn~\eqref{eq:action}, phần thưởng $=$ số quân ăn~\eqref{eq:reward},
hàm chuyển~\eqref{eq:t3}, kết thúc~\eqref{eq:terminal}, xuất phát từ $s_0$.
Mỗi phe tối đa hoá $\sum_t r_c(s_t,a_t)+(\pm W)$---ăn càng nhiều quân càng
tốt, và trên hết là thắng.

\paragraph{Hai cách đọc tác nhân.} Có thể coi \emph{hai người chơi} là
tác nhân ra quyết định (chuẩn lý thuyết trò chơi), khi đó 28 quân là 28
\emph{thực thể bền vững}: loại không đổi, vị trí và phe thay đổi, không bao
giờ biến mất. Hoặc đọc theo kiểu \emph{28 quân là 28 tác nhân chia hai
đội}---bị ăn nghĩa là một quân \emph{đào tẩu} sang đội đối phương; đội
không mất tác nhân nào, chỉ mất lòng trung thành của quân ấy, và tổng hai
đội vẫn là 28. Hai cách đọc mô tả cùng một $\mathcal G$; cách thứ hai nhấn
mạnh đúng tinh thần ``luôn 28 quân, ăn là đổi phe''.
